\section{Introduction}

An advertiser with a limited budget has to decide whom to show an ad to. The
habitual answer is to rank users by their probability of converting and
treat from the top. That ranking answers a \emph{predictive} question: who is
likely to buy? The decision the advertiser actually faces is \emph{causal}:
whose behaviour does the ad change? A user with a 40\,\% chance of buying
anyway is a poor target if the ad moves that chance to 41\,\%, and a user
with a 2\,\% chance who moves to 6\,\% is a good one. Ranking by response and
ranking by effect coincide only if the ad's lift is constant, or
proportional to baseline propensity in a way that preserves the ordering.

This paper asks a narrow empirical question on a public benchmark: can we
identify users who convert \emph{because of} an ad, rather than users who
would have converted anyway, and does targeting them beat targeting the
likeliest converters?

We work on the Criteo Uplift Prediction dataset \cite{diemert2021}, 13,979,592
users pooled from several advertisers' randomized incrementality tests, with
a binary treatment (ad eligibility), two outcomes (site visit and
conversion), and 12 anonymized features. Rather than fit a single uplift
model, we run the whole classical statistical-ML workflow on one fixed
train/validation/test protocol, so that every stage is checked against the
stages before it.

Our findings, each with a confidence interval in Section~\ref{sec:experiments}:
\begin{itemize}
  \item \textbf{Randomization is imperfect and it matters.} Treatment is
        weakly predictable from features, and the imbalance correlates with
        outcomes. Covariate-adjusted ATE estimators agree with each other
        and sit well below the raw difference.
  \item \textbf{Assignment is not exposure.} Only 3.6\,\% of eligible users
        were shown an ad. The naive exposed-versus-control contrast is
        confounded by user activity; an instrumental-variable estimate is
        not.
  \item \textbf{Uplift targeting wins for visits, not for conversions.} At a
        5\,\% budget the X-learner finds about 60\,\% more incremental visits
        than a response model. For conversions, the response model is never
        beaten, and the causal forest is significantly worse on Qini.
\end{itemize}

\paragraph{Contributions.} This paper makes three contributions:
\begin{enumerate}
  \item An end-to-end, reproducible analysis of the Criteo v2.1 benchmark on
        the full 14M rows, with bootstrap intervals on every headline number
        and a test split that is touched once.
  \item A careful treatment of the gap between randomized \emph{assignment}
        and actual \emph{exposure}, and of imbalance in the pooled
        benchmark, which together change the reported effect sizes by up to
        28\,\%.
  \item A direct, budget-matched comparison of uplift and response targeting
        that locates \emph{where} uplift modeling helps (visits at small
        budgets, with at least 3M training rows) and where it does not
        (conversions).
\end{enumerate}

\paragraph{Roadmap.} Section~\ref{sec:related} surveys related work.
Section~\ref{sec:method} describes the data, protocol, and estimators.
Section~\ref{sec:experiments} reports results. Section~\ref{sec:discussion}
interprets them and lists limitations, and Section~\ref{sec:conclusion}
concludes with future work.
